\documentclass[10pt,a4paper]{amsart}
\usepackage[margin=19mm]{geometry}
\usepackage{amsmath,amssymb,mathtools}
\usepackage[scr=rsfs,cal=euler]{mathalfa}
\usepackage{xfrac}
\usepackage[unicode,hidelinks,bookmarks=false]{hyperref}
\newcommand{\ie}{i.e.\ }
\newcommand{\cf}{cf.\ }
\newcommand{\resp}{respectively\ }
\newcommand{\Subsection}{\subsection}
\providecommand{\nobreakdash}{\nobreak\hbox{-}}
\setlength{\emergencystretch}{2em}
\multlinegap0pt
\usepackage{bm}
\usepackage{bbm}
\usepackage{dsfont}
\usepackage{xspace}
\usepackage{cancel}
\usepackage{mathtools}
\usepackage{amsthm}
\makeatletter
\@ifundefined{thm}{\newtheorem{thm}{Thm}}{}
\@ifundefined{lemma}{\newtheorem{lemma}[thm]{Lemma}}{}
\makeatother
\def\C{\mathbb{C}}
\renewcommand{\ker}{\mathrm{Ker}}
\title[Sasakian manifolds and spin-c Killing spinors]{Sasakian manifolds and spin-c Killing spinors}
\author{Alejandro Gil-Garc\'ia, C. S. Shahbazi}
\date{Source: arXiv:2606.11168v1 (CC BY 4.0)}
\begin{document}
\maketitle

\noindent\emph{Source and licence.} Sasakian manifolds and spin-c Killing spinors. Version arXiv:2606.11168v1 (CC BY 4.0), distributed under the Creative Commons Attribution 4.0 International licence, as recorded in the arXiv licence metadata for this exact version and shown on the abstract page https://arxiv.org/abs/2606.11168v1. Full rights notice: licenses/arxiv-2606.11168-CC-BY-4.0.txt.\par\smallskip

\noindent\textit{Editorial scope.} Counted unit: the Lemma whose source label is \texttt{None} in the article (source0.tex lines 883--885, source label \texttt{None}), delivered together with the article's own complete original proof of it (source0.tex lines 887--940, one \texttt{proof} environment of 54 lines). No mathematics is written, repaired or re-derived: statement and proof are the article's own text line for line. Required context retained uncounted: none. Every definition, display and label of those lines is retained unchanged. Omitted from this package and not redistributed: 1--882, 886--886, 941--954. Nothing inside a retained excerpt is omitted or altered: every retained body is delivered exactly as the article's lines are written, so no omission receipt of any kind accompanies this package. Citations: the retained excerpts carry no citation command at all, so no bibliography entry of the article is transcribed and no reference is redistributed. Reference adaptation: none was needed and none was made; every reference inside the delivered unit resolves inside it, so no unresolved reference or citation is produced. Printed slips kept literal: no printed word, symbol, hypothesis or number of the counted statement or of its proof was altered. Layout and class: the article's own class is \texttt{amsart} with packages including geometry, bm, latexsym, amsmath, amstext, amssymb, amsfonts, amscd, array, multirow, amsbsy, mathrsfs, stmaryrd, amsthm; the wrapper is the lane's pinned \texttt{amsart} editorial wrapper at 10pt on A4, which restates the theorem-like environments the retained text needs and adds the article's own macro definitions that the retained text needs (\texttt{\textbackslash C}, \texttt{\textbackslash ker}), with extra wrapper packages (bm, bbm, dsfont, xspace, cancel, mathtools). The delivered numbering is the unit's own. Assets: no retained span contains a figure environment, a graphics-inclusion command, a PDF-page inclusion, a table or a TikZ picture; no image file is copied into this package, the package declares no asset dependency and no image file is redistributed with it. Licence: version arXiv:2606.11168v1 of this article is distributed under the Creative Commons Attribution 4.0 International licence, as recorded in the arXiv licence metadata for this exact version and shown on the abstract page https://arxiv.org/abs/2606.11168v1. A full rights notice is delivered at licenses/arxiv-2606.11168-CC-BY-4.0.txt. Characters: every retained excerpt is pure ASCII, so the delivered engine, XeLaTeX with its default Latin Modern text font, reports no missing character.
\begin{lemma}
The $\mathcal{L}$-valued complex exterior form $\rho\in\Omega^n_\C(M,\mathcal{L})$ satisfies $\nabla^{g,A}_X \rho = i^{n+1}\alpha \ast (X^\flat \wedge \rho)$.
\end{lemma}

\begin{proof}
We have seen that:
\begin{equation}
\label{eq:cov_deriv}
\nabla_{X}^{g,A}\rho = i\varepsilon_n\alpha(\zeta \wedge \iota_{X}\theta) \otimes \mathfrak{l} + (-1)^{\binom{n}{2}}i\alpha\zeta(X)\rho.
\end{equation}

\noindent
Decomposing the vector field $X$ into its horizontal and Reeb components, $X = X_{\mathcal{H}} + \zeta(X)\xi$, we have the dual decomposition $X^{\flat} = X_{\mathcal{H}}^{\flat} + \zeta(X)\zeta$. Applying the Hodge star operator linearly, we get:
\begin{equation*}
*(X^{\flat} \wedge \rho) = *(X_{\mathcal{H}}^{\flat} \wedge \theta)\otimes \mathfrak{l} + \zeta(X)*(\zeta \wedge \theta)\otimes \mathfrak{l}.
\end{equation*}

\noindent
For the horizontal component, on a $(2n+1)$-dimensional contact metric manifold with volume form $\nu=\nu_{\mathcal{H}}\wedge\zeta$, the Hodge star operator reduces to the $2n$-dimensional Hodge star operator on the contact distribution $\mathcal{H}=\ker(\zeta)$. Using that:
\begin{equation*}
*_{\mathcal{H}}\theta=(-1)^{\binom{n+1}{2}}i^n\theta
\end{equation*}

\noindent
for the $(n,0)$-form $\theta$ we get:
\begin{equation*}
*(X_{\mathcal{H}}^{\flat} \wedge \theta) = *_{\mathcal{H}}(X_{\mathcal{H}}^{\flat} \wedge \theta) \wedge \zeta = (-1)^n\iota_{X_{\mathcal{H}}}(*_{\mathcal{H}}\theta)\wedge\zeta = (-1)^{\binom{n}{2}}i^n \iota_{X_{\mathcal{H}}} \theta \wedge \zeta.
\end{equation*}

\noindent
Multiplying by the overall constant $i^{n+1}\alpha$, the horizontal term becomes:
\begin{equation*}
i^{n+1}\alpha *(X_{\mathcal{H}}^{\flat} \wedge \theta) = i^{n+1}\alpha (-1)^{\binom{n}{2}}i^n \iota_{X_{\mathcal{H}}} \theta \wedge \zeta = i\varepsilon_n\alpha(\zeta\wedge\iota_X\theta),
\end{equation*}

\noindent
where we have used the fact that $\iota_{X_{\mathcal{H}}}\theta = \iota_{X}\theta$ since $\theta$ is purely horizontal.\medskip

For the vertical component, we have:
\begin{equation*}
*(\zeta\wedge\theta)=*_{\mathcal{H}}\theta=(-1)^{\binom{n+1}{2}}i^n\theta.
\end{equation*}

\noindent
Multiplying it by the constant $i^{n+1}\alpha$ we obtain:
\begin{equation*}
i^{n+1}\alpha\zeta(X)*(\zeta\wedge\theta)=(-1)^{\binom{n+1}{2}}i^ni^{n+1}\alpha\zeta(X)\theta=(-1)^{\binom{n}{2}}i\alpha\zeta(X)\theta.
\end{equation*}

\noindent
Combining the horizontal and vertical terms we get:
\begin{equation*}
i^{n+1}\alpha *(X^{\flat} \wedge \rho) = i\varepsilon_n\alpha(\zeta \wedge \iota_{X}\theta) \otimes \mathfrak{l} + (-1)^{\binom{n}{2}}i\alpha\zeta(X)\rho,
\end{equation*}

\noindent
which matches the expression for $\nabla_{X}^{g,A}\rho$ in Equation \eqref{eq:cov_deriv}, thereby concluding the proof.
\end{proof}

\end{document}
